\section{Methodology}
\label{sec:methodology}

FourShadow is implemented in Vortex~\cite{vortex}, the open-source RISC-V GPU, and its performance is measured there: in processor cycles on SimX, the Vortex cycle-level simulator, and end to end on an FPGA. FourShadow is compared with two NVIDIA GPUs, an RTX 5080 and an H100, for the reasons given in Section~\ref{sec:intro}. Because the NVIDIA GPUs run different toolkits and FourShadow counts simulated cycles, the platforms are compared only by ratios and trends, never by absolute cycles.

\subsection{Hardware Configuration}\label{sec:method_hw}

Table~\ref{tab:config} lists the configuration of each platform; every run uses one FourShadow core and one SM of each GPU. FourShadow's entries are Vortex build options, with the tensor core, warp-group MMA and DXA enabled; only the accumulator count $\mathit{NRC}$ is a kernel parameter; parentheses name the part of Section~\ref{sec:eval} that sweeps a value, and bold marks the default (an issue width cannot exceed NW). We name a FourShadow configuration by its issue width (IW) and its number of warps per core (NW): (IW-1, NW-4) issues from one warp per cycle on a core of four warps; every part that runs four warps per core states it. The configuration matched to the H100 is (IW-4, NW-4) at 32 threads per warp. The H100 has no $\mathit{NRC}$; each $\mathit{NRC}$ is paired with the \texttt{wgmma.m64nNk16} of equal warp-group tile, $N{=}\mathit{NRC}\,\mathit{NT}/(2\,\mathit{tcM}){=}2\,\mathit{NRC}$ (Eq.~\ref{eq:wg_tile}), whose 128 threads each hold $N/2{=}\mathit{NRC}$ FP32 accumulators~\cite{nvidia2026ptxisa}.

\begin{table}[h]
\centering
\caption{Microarchitecture Configuration Comparison between FourShadow and NVIDIA GPUs}
\label{tab:config}
\footnotesize
\setlength{\tabcolsep}{2pt}
\renewcommand{\arraystretch}{0.95}
\begin{tabular}{@{}l l l l@{}}
\toprule
 & \textbf{FourShadow} & \textbf{RTX 5080} & \textbf{H100 SXM} \\
\midrule
Threads/warp (NT) & 4, 8, 16, 32 (\ref{sec:eval_threads}) & 32 & 32 \\
Warps/core (NW) & 4, 16 & 4 per CTA & 4 per CTA \\
Issue width (IW) & 1, 2, 4, 8 (\ref{sec:eval_iw}) & -- & 4 (warpgroup) \\
Accum.\ regs.\ ($\mathit{NRC}$) & 8, 16, 32 (\ref{sec:eval_nrc}) & -- & -- \\
FEDP width & $K$, $2K$ (\ref{sec:eval_fedp2k}) & not visible & not visible \\
Warp-level MMA & TCU, dense, 2:4 & \texttt{mma(.sp)} & \texttt{mma(.sp)} \\
Warp-group MMA & WGMMA, RS, SS & -- & \texttt{wgmma(.sp)} \\
\midrule
Cycles & simulated & \texttt{clock64} & \texttt{clock64} \\
Tech. Node. & ASAP7 & TSMC 4NP & TSMC 4N \\
Clock Freq. & 1.0 GHz & 2.295 GHz & 1.59 GHz \\
Toolchain & \makecell[l]{VOLT~\cite{10.1145/3771775.3786275},\\\texttt{Vortex 3.0}} & \makecell[l]{CUDA 13.0,\\\texttt{sm\_120a}} & \makecell[l]{CUDA 12.4.131,\\\texttt{sm\_90a}} \\


\bottomrule
\end{tabular}
\end{table}

\subsection{Software Configuration}\label{sec:method_sw}

The design-space sweeps run $128\times128\times K$ SGEMM kernels, where $K\in\{128,256,512,1024\}$, change one architectural parameter while holding the rest of the configuration fixed, and compute speedups from total kernel cycles against the baseline configuration each part names in the (IW-N, NW-M) form; the figures show these kernels. \phantomsection\label{sec:large_method}We repeat these sweeps, except the descriptor load of Section~\ref{sec:eval_desc}, at $256{\times}256{\times}K$ and $512{\times}512{\times}K$ with $K\in\{128,256,512,768,1024,1536\}$, and \phantomsection\label{sec:skinny_method}on skinny GEMMs, tall ($128{\times}n{\times}K$) and wide ($n{\times}128{\times}K$), with $n$ of 16 to 64 and $K$ up to 1536.
% Source (256/512): 00_doc/vortex-wgdxa2k/results/RERUN_256.md and RERUN_512.md header (K in {128, 256, 512, 768, 1024, 1536}; the warp-level halvable-share item F10 also runs 32-2048 at 256 and 128-2048 at 512; items F1-F11 = the skinny item set, no descriptor load), ledgers runs-bfix2-r256-*.jsonl and runs-bfix2-r512-*.jsonl.
% Source (skinny): 00_doc/vortex-wgdxa2k/results/SKINNY.md sections 1 and 5 (S = 128 family; views t16/t32/t64, w16/w64; K 128-1536; items F1-F11 = every 128x128 sweep but the descriptor load), results/skinny.json.
The comparisons with the NVIDIA GPUs run instruction-level kernels and count cycles over an in-kernel region. Warp-level MMA is compared at (IW-1, NW-4) with 8 or 32 threads per warp, because there the FP16 sparse speedups of FourShadow and the RTX 5080, over the 17 shapes measured on both, differ by a median of 9.0\% at 8 threads and 6.2\% at 32 threads, less than at issue widths of 2 and 4.
